The standardized citation metrics author database of Ioannidis and colleagues, widely known as the Stanford top-2\% list, ranks researchers by a composite of six cumulative citation indicators computed from Scopus. We describe an open alternative built from OpenAlex and from the bipartite spectral ranking of sources and institutions. Each work receives a publication score, from the size-independent influence of its source and institutions, and a citer score, the mean influence of the sources of the works that cite it. A researcher's score is the mean over their papers, with full or fractional counting and with or without self-citations. We compute scores for the \NOaEligible{} OpenAlex authors who meet the Stanford volume baseline of five papers and match \NMatchCareerPct{}\% of the Stanford career table to OpenAlex authors by name, location and volume. Self-citation shares measured in the two databases agree closely (Spearman \NSelfRho{}). Spectral scores correlate weakly with the Stanford composite among listed researchers (0.27--0.32) and more strongly with Stanford citations per paper (up to \NRhoPubfracNcpp{}), and top-2\% lists drawn up on spectral scores overlap little with the Stanford list, because a mean rewards the typical reception of a researcher's papers rather than their total. Removing self-citations barely changes the citer score and raises it for heavy self-citers, whose self-citations come from less influential sources. An influence-weighted mean therefore offers an author-level indicator that is open, field-normalised and difficult to inflate by self-citation, at the cost of favouring short records unless combined with a measure of volume.
